\chapter{EMBEDDED FIRMWARE AND DIGITAL SIGNAL PROCESSING}
\label{ch:firmware_dsp}

\section{Firmware Architecture and Modular C++ Design}
The VibeGuard firmware is engineered in modern, modular C++ using the Arduino-ESP32 core under the FreeRTOS real-time operating system. The software architecture emphasizes object-oriented encapsulation, deterministic non-blocking execution, zero dynamic heap allocation during steady-state processing, and strict separation of real-time sampling from spectral calculation.

The firmware codebase is divided into five cohesive modules:
\begin{enumerate}
    \item \texttt{adxl345\_spi}: Low-level driver encapsulating 4-wire hardware SPI transactions, register configuration, bitwise masking, and FIFO burst reading.
    \item \texttt{dsp\_pipeline}: Numerical signal processing library executing in-place DC gravity removal, single-axis RMS, triaxial Vector RMS, Hann windowing, and the Cooley-Tukey Radix-2 FFT.
    \item \texttt{persistence\_filter}: Temporal state queue managing circular history buffers for transient glitch rejection.
    \item \texttt{telemetry}: High-speed, non-blocking formatted serial streamer transmitting NMEA-style ASCII telemetry over UART.
    \item \texttt{main}: System executive implementing the 5-state Finite State Machine, hardware timer interrupts, and RGB LED control.
\end{enumerate}

\section{Sensor Bring-up and Device ID Verification}
To prevent the monitoring system from operating with a disconnected, unpowered, or malfunctioning sensor, the firmware implements a strict hardware identity handshake during boot.

Upon power-on reset, the microcontroller initializes the VSPI bus at 5.0\,MHz in SPI Mode~3 ($\mathrm{CPOL}=1, \mathrm{CPHA}=1$). The firmware asserts the active-low Chip Select line (GPIO21) and reads the read-only \texttt{DEVID} register at address \texttt{0x00}. The returned 8-bit value is compared against the factory-defined silicon identity:
\begin{equation}
\texttt{DEVID} = \texttt{0xE5} \quad (1110\,0101_2)
\end{equation}
If the returned value is \texttt{0x00} (bus grounded or sensor unpowered), \texttt{0xFF} (floating MISO bus), or any mismatched value, the bootloader aborts initialization immediately. The system enters \texttt{STATE\_FAULT}, lights the RGB LED solid White, and transmits an error diagnostic packet over serial.

Once \texttt{DEVID} is verified, the configuration sequence writes the following register values:
\begin{itemize}
    \item \textbf{Rate Register (\texttt{BW\_RATE}, \texttt{0x2C} $\leftarrow$ \texttt{0x0D}):} Configures the internal output data rate to $800$\,Hz. Bit 4 (\texttt{LOW\_POWER}) is cleared to maintain full-bandwidth low-noise analog performance.
    \item \textbf{Format Register (\texttt{DATA\_FORMAT}, \texttt{0x31} $\leftarrow$ \texttt{0x08}):} Asserts Bit 3 (\texttt{FULL\_RES}) for constant $3.9\,\text{mg/LSB}$ sensitivity and clears Bits 1--0 to select the $\pm 4g$ dynamic range.
    \item \textbf{FIFO Control (\texttt{FIFO\_CTL}, \texttt{0x38} $\leftarrow$ \texttt{0x80}):} Configures the FIFO buffer in Stream Mode, ensuring continuous rolling buffering of 32 triaxial samples.
    \item \textbf{Power Control (\texttt{POWER\_CTL}, \texttt{0x2D} $\leftarrow$ \texttt{0x08}):} Sets Bit 3 (\texttt{Measure}), transitioning the internal state machine from standby to active continuous measurement.
\end{itemize}

\section{Real-Time Acquisition Loop and Ping-Pong Buffering}
At an 800\,Hz sampling rate, a fresh 6-byte triaxial sample arrives every $1250\,\mu\text{s}$. If the microcontroller attempts to process an analysis window within the same memory buffer being updated by the sensor, race conditions and data corruption inevitably occur.

To guarantee zero sample dropouts, VibeGuard implements a \textbf{Ping-Pong (Double Buffering)} memory architecture. Two static, contiguous memory buffers (\texttt{Buffer\_A} and \texttt{Buffer\_B}) are allocated in internal SRAM:
\begin{itemize}
    \item Each buffer contains three 256-element arrays of 16-bit signed integers ($256 \times 2\,\text{bytes} \times 3\,\text{axes} = 1,536\,\text{bytes}$ per buffer).
    \item While the high-priority sampling routine fills \texttt{Buffer\_A} via SPI at 800\,Hz, the background DSP task executes DC removal, Vector RMS, and FFT analysis on \texttt{Buffer\_B}.
    \item Upon the arrival of the 256th sample, an atomic pointer swap occurs in less than $0.2\,\mu\text{s}$. \texttt{Buffer\_B} becomes the active acquisition target, while \texttt{Buffer\_A} is submitted to the DSP pipeline.
\end{itemize}
This decoupled architecture guarantees deterministic real-time data acquisition with zero dropped blocks, regardless of background DSP execution time.

\section{Static Component and DC Gravity Removal}
Raw digital acceleration outputs ($x_i, y_i, z_i$) contain both dynamic mechanical oscillations and static gravitational projections. To eliminate orientation dependency, the dynamic signal is extracted via two-pass in-place DC cancellation.

In the first pass across the active 256-sample buffer, the arithmetic mean is accumulated:
\begin{equation}
\bar{x} = \frac{1}{N} \sum_{i=0}^{N-1} x_i, \quad \bar{y} = \frac{1}{N} \sum_{i=0}^{N-1} y_i, \quad \bar{z} = \frac{1}{N} \sum_{i=0}^{N-1} z_i
\end{equation}
In the second pass, the static mean is subtracted in-place, and the resulting integer difference is converted to engineering units of $g$ ($1g = 9.80665\,\text{m/s}^2$) using the constant scale factor $S = 0.0039\,g/\text{LSB}$:
\begin{equation}
x'_i = (x_i - \bar{x}) \cdot S, \quad y'_i = (y_i - \bar{y}) \cdot S, \quad z'_i = (z_i - \bar{z}) \cdot S
\end{equation}
Because the operation occurs entirely within existing RAM buffers, memory allocation overhead is zero.

\section{Vector RMS Calculation and Anomaly Metric}
Root-Mean-Square (RMS) vibration magnitude is computed independently for each orthogonal measurement axis:
\begin{equation}
\mathrm{RMS}_x = \sqrt{\frac{1}{N} \sum_{i=0}^{N-1} (x'_i)^2}, \quad \mathrm{RMS}_y = \sqrt{\frac{1}{N} \sum_{i=0}^{N-1} (y'_i)^2}, \quad \mathrm{RMS}_z = \sqrt{\frac{1}{N} \sum_{i=0}^{N-1} (z'_i)^2}
\end{equation}
The composite \textbf{Vector RMS} ($a_{\mathrm{vec}}$) is synthesized as the Euclidean norm of the orthogonal components:
\begin{equation}
a_{\mathrm{vec}} = \sqrt{\mathrm{RMS}_x^2 + \mathrm{RMS}_y^2 + \mathrm{RMS}_z^2}
\label{eq:firmware_vec_rms}
\end{equation}
Vector RMS provides two decisive advantages for embedded condition monitoring:
\begin{enumerate}
    \item \textbf{Rotational Invariance:} If the sensor is mounted at an arbitrary angle relative to the machine frame, single-axis RMS values fluctuate depending on alignment. The Vector RMS sum remains mathematically invariant to rigid sensor rotation.
    \item \textbf{Multi-Planar Sensitivity:} Rotor unbalance generates dynamic forces in the radial ($X\text{--}Y$) plane, while shaft misalignment and foundation rocking generate severe axial ($Z$) oscillations. Vector RMS captures all vibrational modes simultaneously in a single scalar feature.
\end{enumerate}

\section{$K$-of-$M$ Persistence Filter for Glitch Rejection}
Directly triggering alarms based on single-window threshold crossings produces unacceptable false-alarm rates in industrial environments. Benign events---such as an operator resting a tool on the table, nearby door closures, or electrical line transients---can cause a single 320\,ms window to exceed the threshold.

To solve this, VibeGuard implements a temporal \textbf{$K$-of-$M$ Persistence Filter}.
\begin{itemize}
    \item A circular history buffer of length $M = 5$ maintains the threshold violation status of the last five consecutive analysis windows.
    \item For each completed window $j$, a binary violation flag $v_j$ is evaluated:
    \begin{equation}
    v_j = \begin{cases} 
    1 & \text{if } a_{\mathrm{vec}} > \theta_{\mathrm{RMS}} \\
    0 & \text{if } a_{\mathrm{vec}} \le \theta_{\mathrm{RMS}}
    \end{cases}
    \end{equation}
    \item The total active violation count is summed across the history queue:
    \begin{equation}
    V = \sum_{j=1}^{M} v_j
    \end{equation}
    \item An alarm state transition (\texttt{STATE\_ABNORMAL}) is asserted if and only if:
    \begin{equation}
    V \ge K, \quad \text{where } K = 4 \text{ and } M = 5
    \end{equation}
\end{itemize}

\subsection{Filter Characteristics and Latency Analysis}
An isolated mechanical jolt produces at most one or two elevated windows ($V \le 2 < 4$), which the persistence filter completely ignores. The system remains in \texttt{STATE\_NORMAL} (Green LED), eliminating false alarms.

Conversely, a true continuous mechanical defect (such as rotor unbalance) causes sustained threshold violations. The alarm condition is asserted on the 4th consecutive anomalous window ($K = 4$). The physical alert detection latency $\tau_{\mathrm{alert}}$ is:
\begin{equation}
\tau_{\mathrm{alert}} = K \cdot T_{\mathrm{window}} = 4 \times 320\,\text{ms} = 1.28\,\text{seconds}
\end{equation}
A detection latency of $1.28\,\text{s}$ comfortably satisfies the project requirement ($\tau_{\mathrm{alert}} \le 2.0\,\text{s}$), providing rapid emergency alerting while ensuring absolute immunity to transient bench shocks.

\section{Triggered Cooley-Tukey Radix-2 FFT with Hann Windowing}
To prevent continuous Fourier transforms from degrading acquisition timing, the frequency-domain spectral pipeline is \textbf{triggered on-demand}. The FFT executes only when:
\begin{enumerate}
    \item The persistence filter confirms an abnormal condition ($V \ge K$); or
    \item The system executes baseline calibration; or
    \item An explicit diagnostic command is received over serial.
\end{enumerate}

\subsection{Hann Window Implementation}
Before executing the FFT, the zero-mean signal $x'_n$ of the dominant radial axis is multiplied by the precomputed discrete Hann window coefficients:
\begin{equation}
x''_n = x'_n \cdot w(n) = x'_n \cdot 0.5 \left[ 1 - \cos\left(\frac{2\pi n}{N - 1}\right) \right], \quad n = 0, 1, \dots, 255
\end{equation}
To maximize execution speed, the 256 floating-point values of $w(n)$ are precomputed at compile time and stored in flash memory, eliminating runtime trigonometric calculations.

\subsection{In-Place Radix-2 FFT Execution}
The discrete transform is computed using the Decimation-in-Time (DIT) Cooley-Tukey algorithm. The algorithm executes in three stages:
\begin{enumerate}
    \item \textbf{Bit-Reversal Permutation:} The windowed samples are rearranged into bit-reversed order using an in-place swap routine based on integer bit manipulation.
    \item \textbf{Butterfly Computations:} The transform evaluates $\log_2(256) = 8$ successive stages of dual-node butterfly operations. Complex twiddle factors $W_N^k = e^{-j 2\pi k / N}$ are retrieved from a static 128-entry lookup table.
    \item \textbf{Spectral Magnitude Extraction:} The single-sided magnitude spectrum is calculated for all bins from $k = 0$ to $k = N/2 = 128$:
    \begin{equation}
    |X(k)| = \frac{2}{N} \sqrt{\mathrm{Re}\{X(k)\}^2 + \mathrm{Im}\{X(k)\}^2}, \quad k = 1, 2, \dots, 127
    \end{equation}
    The DC bin ($k = 0$) and Nyquist bin ($k = 128$) are scaled by $1/N$.
\end{enumerate}

\subsection{Harmonic Peak Detection}
The firmware scans the magnitude array $|X(k)|$ across bins $k \in [1, 127]$ to identify the dominant spectral peak:
\begin{equation}
k_{\mathrm{peak}} = \arg\max_{k \ge 1} |X(k)|
\end{equation}
The physical peak vibration frequency $f_{\mathrm{peak}}$ is computed directly:
\begin{equation}
f_{\mathrm{peak}} = k_{\mathrm{peak}} \cdot \Delta f = k_{\mathrm{peak}} \cdot 3.125\,\text{Hz}
\end{equation}
For our 600\,RPM motor, an unbalance anomaly produces a sharp maximum at bin $k = 3$ ($9.375$\,Hz) or bin $k = 4$ ($12.5$\,Hz), matching the theoretical $10.0$\,Hz fundamental.

\section{Memory Footprint and Real-Time Timing Budget}
Table~\ref{tab:memory_budget} summarizes the static RAM consumption and computational timing budget measured on the ESP32 core running at 160\,MHz.

\begin{table}[H]
\centering
\begin{tabular}{|l|c|r|}
\hline
\textbf{Firmware Component} & \textbf{Memory Allocation} & \textbf{Execution Duration} \\
\hline
Ping-Pong Buffers (\texttt{Buffer\_A} + \texttt{Buffer\_B}) & 3,072\,bytes SRAM & $11.2\,\mu\text{s}$ per sample (SPI) \\
In-Place Zero-Mean DC Removal (3 Axes) & In-place (0 extra RAM) & $0.082\,\text{ms}$ per window \\
Triaxial Vector RMS Feature Extraction & 24\,bytes stack & $0.145\,\text{ms}$ per window \\
$K$-of-$M$ Persistence Filter Queue & 10\,bytes SRAM & $0.004\,\text{ms}$ per window \\
Hann Windowing (Lookup Table in Flash) & 1,024\,bytes Flash & $0.065\,\text{ms}$ per FFT trigger \\
256-Point Radix-2 Complex FFT & 2,048\,bytes SRAM & $1.020\,\text{ms}$ per FFT trigger \\
Formatted USB Serial Telemetry Output & 128\,bytes buffer & $0.045\,\text{ms}$ (DMA transfer) \\
\hline
\textbf{Total Routine Time (Time-Domain Only)} & \textbf{3,106\,bytes RAM} & \textbf{0.276\,ms per 320\,ms window} \\
\textbf{Total Peak Time (With Triggered FFT)} & \textbf{5,154\,bytes RAM} & \textbf{1.361\,ms per 320\,ms window} \\
\hline
\end{tabular}
\caption{VibeGuard firmware memory footprint and processor timing budget}
\label{tab:memory_budget}
\end{table}

As demonstrated in Table~\ref{tab:memory_budget}, even during windows where a complete 256-point FFT is triggered, total DSP execution consumes only $1.361\,\text{ms}$ out of the available $320.0\,\text{ms}$ window duration. This represents a CPU utilization of less than \textbf{0.43\%}, leaving over \textbf{99.5\%} of microcontroller processing headroom completely unencumbered.

\section{Summary}
This chapter detailed the embedded firmware architecture and digital signal processing algorithms implemented in VibeGuard. Modular C++ drivers established reliable SPI Mode~3 communication, verified sensor identity via register \texttt{DEVID} (\texttt{0xE5}), and eliminated race conditions using a double-buffered Ping-Pong memory scheme. In-place zero-mean cancellation removed static gravity vectors without dynamic RAM allocation. Triaxial Vector RMS synthesized orthogonal vibration into an orientation-invariant scalar metric. A $K$-of-$M$ persistence filter ($K=4, M=5$) rejected transient bench jolts while providing an alert latency of $1.28$\,seconds. Finally, a triggered 256-point Cooley-Tukey FFT with Hann windowing isolated mechanical harmonic peaks with 3.125\,Hz resolution, executing in just $1.02$\,ms with an overall CPU load under 0.5\%.
